% NichidaiMasterExam_R8_2nd_1
\begin{tcolorbox} %%R8_2nd_1
	$\fbox{1}$
\begin{enumerate}
    \item $V, W$をｔ複素線形空間とし、$V, W$の零ベクトルをそれぞれ$\bm{0}_V, \bm{0}_W$とする。$V$から$W$への線形写像$T\colon V\to W$に対して、$T$の像を$\Im{(T)}$とし、$T$の核を$\Ker{(T)}$とする。$\Im{(T)}$および$\Ker{(T)}$の定義をそれぞれ述べなさい。
    \item $\C$を複素数体とする。線形写像$S\colon \C^3\to \C^4$および$T\colon \C2^\to \C^4$を以下で定める。
\begin{equation*}
    S\left(\begin{pmatrix}x_1& x_2& x_3\end{pmatrix}\right)= 
    \begin{pmatrix}
        3x_1- x_2+ 5x_3\\
        x_1+ x_2+ 3x_3\\
        x_1+ 2x_3\\
        x_1- 3x_2- x_3
    \end{pmatrix},\ 
    T\left(\begin{pmatrix}x_1\\ x_2\end{pmatrix}\right)= 
    \begin{pmatrix}
        x_1- 3x_2\\
        3x_1+ x_2\\
        x_1- 2x_2\\
        -5x_1- x_2
    \end{pmatrix}
\end{equation*}
\begin{enumerate_alpha}
    \item $\Ker{(S)}= \begin{Bmatrix}\alpha\begin{pmatrix}-2\\ -1\\ 1\end{pmatrix}; \alpha\in \C\end{Bmatrix}$を示しなさい。
    \item $\Im{(S)}$と$\Im{(S)}$の和空間を$W= \Im{(S)}+ \Ker{(T)}$とおく。$W$の次元を求めなさい。またその根拠も述べなさい。

\end{enumerate}
\end{tcolorbox}

\begin{souanswer} %R8_2nd_1
\begin{enumerate}
    \item 
\begin{tcolorbox}[title= 像・核の定義]
\begin{itemize}
    \item 像($\img$)
\begin{equation*}
    \img{(T)}= \{T(v)\in W; v\in W\}
\end{equation*}
    \item 核($\ker$)
\begin{equation*}
    \ker{(T)}= \{v\in V; T(v)= \bm{0}_W\}
\end{equation*}
\end{itemize}
\end{tcolorbox}
    \item 
\begin{enumerate_alpha}
    \item $\bm{x}= \begin{pmatrix}x_1\\ x_2\\ x_3\end{pmatrix}\in \ker{(S)}\iff S(\bm{x})= \bm{0}$であるから、次の連立方程式を解く。
\begin{equation*}
\begin{pmatrix}
    3& -1& 5\\
    1& 1& 3\\
    1& 0& 2\\
    1& -3& -1
\end{pmatrix}\begin{pmatrix}x_1\\ x_2\\ x_3\end{pmatrix}
    = \begin{pmatrix}0\\ 0\\ 0\\ 0\end{pmatrix}
\end{equation*}
    係数行列を基本変形する。
\begin{equation*}
    \begin{pmatrix}
    3& -1& 5\\
    1& 1& 3\\
    1& 0& 2\\
    1& -3& -1
\end{pmatrix}\Rightarrow
\begin{pmatrix}
    1& 0& 2\\
    0& 1& 1\\
    0& 0 &0\\
    0& 0 &0\\
\end{pmatrix}
\end{equation*}
    簡約化された行列より、連立方程式は以下と同値。
\begin{equation*}
\begin{cases}
    x_1+ 2x_3= 0\\
    x_2+ x_3= 0
\end{cases}
    \iff
\begin{cases}
    x_1= -2x_3\\
    x_2= -x_3
\end{cases}
\end{equation*}
    $x_3= \alpha\in \C$とおくと、$x_3= \alpha\beign{pmatrix}-2\\ -1\\ 1\end{pmatrix}$。\\
    したがって$\ker{(S)}= \{\alpha\beign{pmatrix}-2\\ -1\\ 1\end{pmatrix}; \alpha\in \C\}$が示された。
    
    \item $\dim(\img(S))= 1$であり、次元定理より$\dim(\img(S))= 3- 1= 2$。\\
    $S$の表現行列の第一・二列は線形独立であるため、$\img(S)$の基底として次が取れる。
\begin{equation*}
    v_1= \begin{pmatrix}3\\ 1\\ 1\\ 1\end{pmatrix},\ 
    v_2= \begin{pmatrix}-1\\ 1\\ 0\\ 3\end{pmatrix}
\end{equation*}
    $T$の表現行列の列ベクトルを$w_1= \begin{pmatrix}1\\ 3\\ 1\\ -5\end{pmatrix},\ w_2= \begin{pmatrix}-3\\ 1\\ -2\\ -1\end{pmatrix}$とおくと、$w_1, w_2$は明らかに1次元であるため、$\img(T)= \span\{w_1, w_2\}$。\\
    $W= \span\{v_1, v_2, w_1, w_2\}$であるため、これらのベクトルを並べた行列の階数を求める。
\begin{equation*}
    A= 
\begin{pmatrix}
    3& -1& 1& -3\\
    1& 1& 3& 1\\
    1& 0& 1& -2\\
    1& 3& -5& -1
\end{pmatrix}
\end{equation*}
    行基本変形を行うと、
\begin{equation*}
    A= 
\begin{pmatrix}
    1& 0& 1& -2\\
    0& 1& 2& 3\\
    0& 0& 0& 6\\
    0& 0& 0& 0
\end{pmatrix}
\end{equation*}
    非零行は3行であるため、$\rank{A}= 3$である。\\
    したがって$W$の次元は3である。
\end{enumerate_alpha} 
\end{enumerate}
\end{souanswer}
